\section*{Hoofdstuk \ref{ch:b3:banach} --- Banachruimten en de
fundamentele stellingen}

\begin{solution}{exo:b3:banach:1}
(a) $\norm{Sx}_2 = \norm x_2$: $S$ is een isometrie, dus
$\vertiii S = 1$. Voor de achterwaartse verschuiving is
$\norm{S^*x}_2^2 = \sum_{n \geq 2}\abs{x_n}^2 \leq \norm x_2^2$,
met gelijkheid voor $x = e_2$: $\vertiii{S^*} = 1$.

(b) $\norm{M_ax}_2^2 = \sum\abs{a_n}^2\abs{x_n}^2 \leq \norm
a_\infty^2\norm x_2^2$; en toetsen op $x = e_n$ geeft
$\vertiii{M_a} \geq \abs{a_n}$ voor elke $n$: $\vertiii{M_a} =
\norm a_\infty$.

(c) $\abs{\Lambda(f)} \leq \int_0^1\abs f \leq \norm f_\infty$,
dus $\norm\Lambda \leq 1$. Zij voor $\varepsilon > 0$ de functie
$f_\varepsilon$ gelijk aan $1$ op $[0, \frac12 - \varepsilon]$,
aan $-1$ op $[\frac12 + \varepsilon, 1]$, en affien daartussen:
dan is $\norm{f_\varepsilon}_\infty = 1$ en
$\Lambda(f_\varepsilon) \geq 1 - 2\varepsilon$, dus $\norm\Lambda
= 1$. Niet bereikt: uit $\Lambda(f) = 1$ met $\norm f_\infty \leq
1$ volgt $\int_0^{1/2}f = \frac12$ en $\int_{1/2}^1 f =
-\frac12$, dat wil zeggen (\omterm{def:b3:topology:continuity}{continuïteit} en $\abs f \leq 1$) $f
\equiv 1$ op $[0, \frac12]$ en $f \equiv -1$ op $[\frac12, 1]$:
tegenspraak in $\frac12$.
\end{solution}

\begin{solution}{exo:b3:banach:2}
Schrijf $S = T\bigl(I - T^{-1}(T - S)\bigr)$ met
$\vertiii{T^{-1}(T-S)} \leq \vertiii{T^{-1}}\,\vertiii{T - S} =
\theta < 1$: volgens \cref{prop:b3:banach:neumann} is $S$
inverteerbaar met $S^{-1} = \sum_{n\geq0}\bigl(T^{-1}(T -
S)\bigr)^nT^{-1}$, waaruit
\[
\vertiii{S^{-1} - T^{-1}} \leq
\sum_{n\geq1}\theta^n\,\vertiii{T^{-1}}
= \frac{\vertiii{T^{-1}}^2\,\vertiii{T - S}}{1 - \theta}.
\]
Voor het lineaire stelsel: $x_T = T^{-1}b$ en $x_S = S^{-1}b$
verschillen hoogstens die grens maal $\norm b$ --- kleine
verstoringen van een inverteerbaar stelsel blijven eenduidig
\omterm{def:b3:groups:derived}{oplosbaar}, met lipschitz-afhankelijkheid van de oplossing van de
operator.
\end{solution}

\begin{solution}{exo:b3:banach:3}
(a) $\ell^1$: zij $(x^{(k)})$ Cauchy. Elke coördinaat is Cauchy
($\abs{x^{(k)}_n - x^{(l)}_n} \leq \norm{x^{(k)} - x^{(l)}}_1$);
stel $x_n = \lim_kx^{(k)}_n$. Bij gegeven $\varepsilon$ is voor
$k, l \geq K$: $\sum_{n \leq N}\abs{x^{(k)}_n - x^{(l)}_n} \leq
\varepsilon$ voor elke $N$; laat $l \to \infty$ en dan $N \to
\infty$: $\norm{x^{(k)} - x}_1 \leq \varepsilon$, en $x = x^{(k)}
- (x^{(k)} - x) \in \ell^1$. $\ell^\infty$: Cauchy voor
$\norm\cdot_\infty$ is uniform Cauchy, en convergeert dus uniform
naar een begrensde rij. $c_0$ is gesloten in $\ell^\infty$:
convergeert $x^{(k)} \to x$ uniform met $x^{(k)}_n \to_n 0$, dan
geeft $\abs{x_n} \leq \norm{x - x^{(k)}}_\infty +
\abs{x^{(k)}_n}$ dat $\limsup_n\abs{x_n} \leq \varepsilon$: dus
$x \in c_0$; en een gesloten deelruimte van een banachruimte is
een banachruimte. Eindige rijen: hun \omterm{def:b3:topology:interior}{afsluiting} bevat elke $x \in
c_0$ (de afknottingen convergeren: $\sup_{n>N}\abs{x_n} \to 0$)
en ligt binnen het gesloten $c_0$.

(b) Neem wegens de homogeniteit $\norm x_p = 1$: dan is
$\abs{x_n} \leq 1$ voor alle $n$, dus $\abs{x_n}^q \leq
\abs{x_n}^p$ en $\norm x_q \leq 1 = \norm x_p$; voor $q = \infty$
geldt $\abs{x_n} \leq \norm x_p$ rechtstreeks. Striktheid: $x_n =
n^{-\alpha}$ met $\frac1q < \alpha \leq \frac1p$ ligt in $\ell^q
\setminus \ell^p$ (riemannreeksen).
\end{solution}

\begin{solution}{exo:b3:banach:4}
($\leq$) Is $\norm f \leq 1$ en $f\restriction_F = 0$, dan is voor
elke $y \in F$: $\abs{f(x)} = \abs{f(x - y)} \leq \norm{x - y}$;
neem het infimum. ($\geq$, en bereikt)
\cref{cor:b3:banach:hbnormed}(3) levert een $f$ met
$f\restriction_F = 0$, $\norm f \leq 1$ en $f(x) = d(x, F)$: het
supremum is dus een maximum. Gevolg: $F \subseteq \bigcap\{\ker f
: f\restriction_F = 0\}$ is triviaal, en een punt $x \notin F$
wordt door de bovengenoemde functionaal uit de doorsnede
uitgesloten ($f(x) = d(x,F) > 0$, want $F$ is gesloten).
\end{solution}

\begin{solution}{exo:b3:banach:5}
Voor $y \in \ell^\infty$ is $\abs{\Lambda_y(x)} = \abs{\sum
x_ny_n} \leq \norm y_\infty\norm x_1$, dus $\norm{\Lambda_y} \leq
\norm y_\infty$; en toetsen op $e_n$ geeft $\abs{y_n} =
\abs{\Lambda_y(e_n)} \leq \norm{\Lambda_y}$: gelijkheid. Stel
omgekeerd $y_n = \Lambda(e_n)$ voor $\Lambda \in (\ell^1)'$: dan
is $\abs{y_n} \leq \norm\Lambda$, dus $y \in \ell^\infty$; en
$\Lambda$ en $\Lambda_y$ vallen samen op de eindige rijen, die
\omterm{def:b3:topology:interior}{dicht} liggen in $\ell^1$ ($\norm{x - \sum_{n\leq N}x_ne_n}_1 =
\sum_{n>N}\abs{x_n} \to 0$): dus $\Lambda = \Lambda_y$. De
afbeelding $y \mapsto \Lambda_y$ is lineair, isometrisch en
surjectief. Voor $(\ell^\infty)'$ levert hetzelfde begin een rij
$y_n = \Lambda(e_n)$, maar de eindige rijen liggen \emph{niet}
\omterm{def:b3:topology:interior}{dicht} in $\ell^\infty$ (de constante rij $\mathbf 1$ ligt op
afstand $1$ van alle), zodat $\Lambda$ niet door de $y_n$ wordt
vastgelegd --- en inderdaad is $(\ell^\infty)' \neq \ell^1$
(\cref{pb:b3:banach:1}).
\end{solution}

\begin{solution}{exo:b3:banach:6}
Voor elke vaste $x$ is $y \mapsto B(x, y)$ \omterm{def:b3:topology:continuity}{continu} en lineair:
$\sup_{\norm y \leq 1}\norm{B(x,y)} < \infty$. Dus is de familie
$\{B(\cdot, y) : \norm y \leq 1\} \subseteq \mathcal L(E, G)$
(elk lid \omterm{def:b3:topology:continuity}{continu}, wegens de \omterm{def:b3:topology:continuity}{continuïteit} in $x$) puntsgewijs
begrensd op de banachruimte $E$: Banach--Steinhaus
(\cref{thm:b3:banach:banachsteinhaus}) levert een $C$ met
$\norm{B(x,y)} \leq C\norm x$ voor alle $\norm y \leq 1$; en de
homogeniteit in $y$ maakt het af: $\norm{B(x,y)} \leq C\norm
x\norm y$.
\end{solution}

\begin{solution}{exo:b3:banach:7}
(a) $\norm f_1 \leq \norm f_\infty$: de identiteit is begrensd en
bijectief. Haar inverse is onbegrensd: $f_n(x) = x^n$ heeft
$\norm{f_n}_1 = \frac1{n+1} \to 0$ maar $\norm{f_n}_\infty = 1$.
Geen tegenspraak met \cref{cor:b3:banach:equivnorms}: $(\mathcal
C(\intcc01), \norm\cdot_1)$ is \emph{niet} \omterm{def:b3:complete:complete}{volledig}
(\cref{exo:b3:complete:1}); het gevolg vergt \omterm{def:b3:complete:complete}{volledigheid} aan
beide kanten.

(b) Op de ruimte $E_0$ van de eindige rijen (\omterm{def:b3:topology:interior}{dicht} in $\ell^2$)
is $\varphi(x) = \sum_n n\,x_n$ lineair en onbegrensd
($\varphi(e_n) = n$ met $\norm{e_n}_2 = 1$). De
gesloten-grafiekstelling is niet van toepassing: $E_0$ is niet
\omterm{def:b3:complete:complete}{volledig} --- en $\varphi$ heeft geen \omterm{def:b3:topology:continuity}{continue} voortzetting tot
$\ell^2$, wat illustreert dat dichtheid zonder uniforme
\omterm{def:b3:topology:continuity}{continuïteit} machteloos is
(\cref{thm:b3:complete:extension}).
\end{solution}

\begin{solution}{exo:b3:banach:8}
We gaan de hypothese van de gesloten grafiek na. Zij $x_k \to x$
en $Tx_k \to z$ in $\ell^2$. Dan is voor elke $y$:
\[
\langle z, y\rangle = \lim_k\langle Tx_k, y\rangle
= \lim_k \langle x_k, Ty\rangle = \langle x, Ty\rangle
= \langle Tx, y\rangle,
\]
met de \omterm{def:b3:topology:continuity}{continuïteit} van het inproduct in elke ingang
(Cauchy--Schwarz) en tweemaal de symmetrie. Dus staat $z - Tx$
loodrecht op elke $y$, en in het bijzonder op zichzelf: $z = Tx$.
De grafiek is gesloten en $\ell^2$ is een banachruimte: dus is $T$
begrensd (\cref{thm:b3:banach:closedgraph}). Een symmetrische
operator die op \emph{heel} $\ell^2$ gedefinieerd is, is dus
automatisch begrensd; werkelijk onbegrensde symmetrische
operatoren (plaats, impuls, hamiltonianen) moeten op echte \omterm{def:b3:topology:interior}{dichte}
deelruimten leven.
\end{solution}

\begin{solution}{exo:b3:banach:9}
Ten eerste is $\norm{\Lambda_n} = \sum_i\abs{w_{i,n}}$: $\leq$ is
de driehoeksongelijkheid; $\geq$ door te toetsen op een
stuksgewijs lineaire $f$ met $\norm f_\infty \leq 1$ en
$f(x_{i,n}) = \operatorname{sign}(w_{i,n})$ (interpoleer lineair
tussen de eindig veel knooppunten; waar knooppunten samenvallen,
vallen de tekens samen).

($\Rightarrow$) Puntsgewijze convergentie in elke $f$ geeft (i),
en puntsgewijze begrensdheid, dus levert Banach--Steinhaus
(\cref{thm:b3:banach:banachsteinhaus}) op de banachruimte
$\mathcal C(\intcc01)$ de uitspraak (ii).

($\Leftarrow$) Zij $M = \sup_n\norm{\Lambda_n} + 1$. Kies bij
gegeven $f$ en $\varepsilon$ een veelterm $P$ met $\norm{f -
P}_\infty < \varepsilon/(2M)$
(\cref{cor:b3:complete:weierstrass}); dan is
\[
\Bigl|\Lambda_n f - \int_0^1 f\Bigr|
\leq \abs{\Lambda_n(f - P)} + \Bigl|\Lambda_nP -
\int_0^1P\Bigr| + \Bigl|\int_0^1(P - f)\Bigr|
\leq \varepsilon + \Bigl|\Lambda_nP - \int_0^1P\Bigr| \to
\varepsilon .
\]
Positieve gewichten met exactheid op de constanten:
$\sum_i\abs{w_{i,n}} = \sum_iw_{i,n} = \Lambda_n(\mathbf 1) =
\int_0^1 1 = 1$ voor regels die exact zijn op de constanten ---
(ii) geldt dan met constante $1$.
\end{solution}

\begin{solution}{exo:b3:banach:10}
De lineariteit van $J$ is formeel; en $\norm{J(x)} = \sup_{\norm
f \leq 1}\abs{f(x)} = \norm x$ volgens
\cref{cor:b3:banach:hbnormed}(2). Is $\dim E = n$, dan is $\dim
E' = n$ (een basis geeft coördinaatfunctionalen), dus $\dim E'' =
n$, en is de injectieve (isometrische) $J$ surjectief.
Separabiliteit: zij $(f_n)$ \omterm{def:b3:topology:interior}{dicht} in $E'$ en kies $\norm{x_n} = 1$
met $\abs{f_n(x_n)} \geq \frac12\norm{f_n}$. Stel $F =
\overline{\operatorname{Vect}}(x_n)$; was $F \neq E$, neem dan een
$g \in E'$ met $g \neq 0$ die op $F$ verdwijnt
(\cref{cor:b3:banach:hbnormed}(3)) en kies $f_{n_k} \to g$:
\[
\norm{f_{n_k} - g} \geq \abs{(f_{n_k} - g)(x_{n_k})}
= \abs{f_{n_k}(x_{n_k})} \geq \tfrac12\norm{f_{n_k}}
\geq \tfrac12\bigl(\norm g - \norm{f_{n_k} - g}\bigr),
\]
dus $\norm{f_{n_k} - g} \geq \frac13\norm g > 0$: tegenspraak.
Bijgevolg is $F = E$, en vormen de combinaties van de $x_n$ met
rationale (of $\Q + \iu\Q$-)coëfficiënten een aftelbare \omterm{def:b3:topology:interior}{dichte
verzameling}.
\end{solution}

\begin{solution}{exo:b3:banach:11}
(a) Goed gedefinieerd: $d(x, F)$ hangt alleen van $\bar x$ af
($x$ over $F$ verschuiven verandert de afstand niet). De
homogeniteit en de driehoeksongelijkheid gaan via het infimum
over van $\norm\cdot$. De scheiding vergt de geslotenheid:
$\norm{\bar x} = 0$ betekent $d(x, F) = 0$, dus $x \in \bar F =
F$, dus $\bar x = 0$. $\vertiii\pi \leq 1$: $\norm{\bar x} \leq
\norm x$. Open bal op open bal: is $\norm{\bar x} < 1$, dan heeft
een zekere representant $\norm{x - y} < 1$; en omgekeerd is
$\pi(B_E(0,1)) \subseteq B_{E/F}(0,1)$ wegens de normongelijkheid
--- dus is $\pi$ open, het modelgeval van de
open-afbeeldingsstelling.

(b) Zij $\sum_k\norm{\bar x_k} < \infty$ en til op met
$\norm{x_k} \leq \norm{\bar x_k} + 2^{-k}$: dan is
$\sum\norm{x_k} < \infty$, dus convergeert $s = \sum_kx_k$ in de
banachruimte $E$ (\cref{exo:b3:complete:1}(b)), en de
\omterm{def:b3:topology:continuity}{continuïteit} van $\pi$ geeft $\sum_k\bar x_k = \bar s$: elke
absoluut convergente reeks van $E/F$ convergeert, en dat is
gelijkwaardig met \omterm{def:b3:complete:complete}{volledigheid} (dezelfde oefening).

(c) De afbeelding $\lambda(x) = \lim_nx_n$ is lineair van $c$
naar $K$ en verdwijnt precies op $c_0$, dus induceert ze een
lineaire bijectie $c/c_0 \to K$. Isometrie: $d(x, c_0) =
\abs{\lambda(x)}$ --- $\leq$: trek van $x$ de rij $x -
\lambda(x)\mathbf 1 \in c_0$ af, zodat $\lambda(x)\mathbf 1$
overblijft, van norm $\abs{\lambda(x)}$; $\geq$: voor $y \in c_0$
is $\norm{x - y}_\infty \geq \limsup_n\abs{x_n - y_n} =
\abs{\lambda(x)}$.
\end{solution}

\begin{solution}{exo:b3:banach:12}
(a) $W = \ker P$ is gesloten (origineel van $0$ onder een \omterm{def:b3:topology:continuity}{continue
afbeelding}); en $V = \operatorname{im}P = \ker(I - P)$ (immers
$Px = x$ precies wanneer $x \in \operatorname{im}P$, met $P^2 =
P$), eveneens gesloten. Elke $x$ splitst als $Px + (x - Px)$ met
$Px \in V$ en $x - Px \in W$, en $V \cap W = 0$ (want $x = Px =
0$): dus $E = V \oplus W$.

(b) Het grafiekargument: zij $x_n \to x$ en $Px_n \to z$. Dan is
$z \in V$ ($V$ gesloten en $Px_n \in V$) en $x_n - Px_n \to x - z
\in W$ ($W$ gesloten). Dus $x = z + (x - z)$ met $z \in V$ en $x -
z \in W$; en de eenduidigheid van de ontbinding geeft $z = Px$. De
grafiek van $P$ is gesloten en $E$ is een banachruimte: dus is $P$
begrensd (\cref{thm:b3:banach:closedgraph}).

(c) (a) en (b) samen vormen de gelijkwaardigheid. In een
hilbertruimte heeft elke gesloten $V$ het gesloten complement
$V^\perp$ (\cref{ch:b3:hilbert}): daar is elke gesloten deelruimte
gecomplementeerd. In algemene banachruimten faalt dat --- $c_0$
heeft in $\ell^\infty$ geen gesloten complement (de stelling van
Phillips, buiten ons gereedschap) --- zodat begrensde projecties
een voorrecht zijn, en de gesloten-grafiekstelling precies de
boekhouding is die meetkundige splitsingen in begrensde
operatoren omzet.
\end{solution}

\begin{solution}{pb:b3:banach:1}
\textbf{1.} Voor $a, b > 0$ geeft de concaviteit van $\log$:
$\log\bigl(\tfrac{a^p}p + \tfrac{b^q}q\bigr) \geq \tfrac1p\log
a^p + \tfrac1q\log b^q = \log(ab)$; exponentieer. (Is $ab = 0$,
dan is de ongelijkheid triviaal.) Gelijkheid precies wanneer $a^p
= b^q$.

\textbf{2.} We mogen $\norm x_p = \norm y_q = 1$ aannemen
(homogeniteit; de nulgevallen zijn triviaal). Dan is
\[
\abs{\langle x, y\rangle} \leq \sum_n\abs{x_n}\abs{y_n}
\leq \sum_n\Bigl(\frac{\abs{x_n}^p}p +
\frac{\abs{y_n}^q}q\Bigr) = \frac1p + \frac1q = 1 =
\norm x_p\norm y_q .
\]
Gelijkheid vergt $\abs{x_n}^p = \abs{y_n}^q$ voor alle $n$ (het
gelijkheidsgeval van Young) plus het uitlijnen van de fasen van
$x_ny_n$.

\textbf{3.} $\abs{x_n + y_n}^p \leq \bigl(\abs{x_n} +
\abs{y_n}\bigr)\abs{x_n + y_n}^{p-1}$; sommeren en Hölder
toepassen ($p$ tegen $q$, met $(p - 1)q = p$) geeft
\[
\norm{x + y}_p^p \leq \bigl(\norm x_p + \norm
y_p\bigr)\,\Bigl(\sum_n\abs{x_n + y_n}^{p}\Bigr)^{1/q}
= \bigl(\norm x_p + \norm y_p\bigr)\,\norm{x+y}_p^{p/q};
\]
is $\norm{x + y}_p \neq 0$ (anders triviaal), deel dan door
$\norm{x+y}_p^{p/q}$ en gebruik $p - \frac pq = 1$. (Dat
$\norm{x+y}_p$ eindig is, komt eerst: $\abs{x_n+y_n}^p \leq
2^p(\abs{x_n}^p + \abs{y_n}^p)$.)

\textbf{4.} \omterm{def:b3:complete:complete}{Volledigheid}: als voor $\ell^1$
(\cref{exo:b3:banach:3}), met coördinaatsgewijze limieten plus de
uniforme staartgrens $\sum_{n\leq N}\abs{x^{(k)}_n -
x^{(l)}_n}^p \leq \varepsilon^p$, waarbij eerst $l$ en dan $N$
naar oneindig gaat. Dichtheid van de eindige rijen: $\norm{x -
\sum_{n\leq N}x_ne_n}_p^p = \sum_{n > N}\abs{x_n}^p \to 0$.

\textbf{5.} Hölder geeft $\abs{\Lambda_y(x)} \leq \norm x_p\norm
y_q$: dus $\norm{\Lambda_y} \leq \norm y_q$. Toetsen: stel
$x^{(N)}_n = \abs{y_n}^{q-1}\overline{\operatorname{sign}}(y_n)$
voor $n \leq N$ en $0$ daarna (met $\operatorname{sign}$ de fase
van modulus $1$, zodat $x_ny_n = \abs{y_n}^q$). Dan is
$\Lambda_y(x^{(N)}) = \sum_{n\leq N}\abs{y_n}^q$ en
$\norm{x^{(N)}}_p = \bigl(\sum_{n\leq N}\abs{y_n}^{q}\bigr)^{1/p}$
(want $(q-1)p = q$), dus
\[
\norm{\Lambda_y} \geq \Bigl(\sum_{n\leq
N}\abs{y_n}^q\Bigr)^{1 - 1/p} \xrightarrow[N\to\infty]{}
\norm y_q .
\]

\textbf{6.} Stel $y_n = \Lambda(e_n)$. Met dezelfde toetsvectoren
is $\sum_{n \leq N}\abs{y_n}^q = \Lambda(x^{(N)}) \leq
\norm\Lambda\,\bigl(\sum_{n\leq N}\abs{y_n}^q\bigr)^{1/p}$,
waaruit $\bigl(\sum_{n\leq N}\abs{y_n}^q\bigr)^{1/q} \leq
\norm\Lambda$ voor elke $N$: dus $y \in \ell^q$ met $\norm y_q
\leq \norm\Lambda$. De functionalen $\Lambda$ en $\Lambda_y$
vallen samen op de \omterm{def:b3:topology:interior}{dichte} eindige rijen (vraag 4): $\Lambda =
\Lambda_y$. Met vraag 5 is $y \mapsto \Lambda_y$ dus een
isometrisch isomorfisme $\ell^q \cong (\ell^p)'$.

\textbf{7.} Zij $\xi \in (\ell^p)''$. Samengesteld met de
isometrie $\ell^q \cong (\ell^p)'$ uit vraag 6 definieert $\xi$
een element van $(\ell^q)'$, dat (vraag 6 met $p$ en $q$
verwisseld) gelijk is aan $\Lambda_z$ voor een unieke $z \in
\ell^p$: voor elke $y \in \ell^q$ is $\xi(\Lambda_y) =
\sum_nz_ny_n$. Anderzijds is $J(z)(\Lambda_y) = \Lambda_y(z) =
\sum_ny_nz_n$: dezelfde waarde. En omdat elk element van
$(\ell^p)'$ een $\Lambda_y$ is, volgt $\xi = J(z)$: $J$ is
surjectief --- $\ell^p$ is \omterm{rem:b3:banach:bidual}{reflexief}.

\textbf{8.} De eindige rijen met rationale ingangen (reëel en
imaginair deel) zijn aftelbaar en liggen \omterm{def:b3:topology:interior}{dicht} in $\ell^p$ ($p <
\infty$) en in $c_0$. In $\ell^\infty$: de familie $\{\mathbf 1_A
: A \subseteq \N\}$ is overaftelbaar met $\norm{\mathbf 1_A -
\mathbf 1_B}_\infty = 1$ voor $A \neq B$; de ballen $B(\mathbf
1_A, \frac12)$ zijn paarsgewijs disjunct, en een \omterm{def:b3:topology:interior}{dichte
verzameling} moet elk ervan snijden: er is dus geen aftelbare
\omterm{def:b3:topology:interior}{dichte verzameling}.

\textbf{9.} Was $\ell^1$ \omterm{rem:b3:banach:bidual}{reflexief}, dan was $(\ell^\infty)' \cong
((\ell^1)')' = J(\ell^1)$ separabel (isometrisch beeld van het
separabele $\ell^1$); en volgens \cref{exo:b3:banach:10} zou de
separabiliteit van de \omterm{def:b3:banach:operator}{duale ruimte} $(\ell^\infty)'$ afdwingen dat
$\ell^\infty$ separabel is --- in strijd met vraag 8. Dus is
$\ell^1$ niet \omterm{rem:b3:banach:bidual}{reflexief} (en is $(\ell^\infty)'$ strikt groter dan
$\ell^1$, zoals vraag 11 concreet maakt).

\textbf{10.} De homogeniteit van $p$ is duidelijk;
subadditiviteit: middelen is lineair, en $\limsup(u_n + v_n) \leq
\limsup u_n + \limsup v_n$. Op $c$: de cesàromiddelen van een
convergente rij convergeren naar haar limiet, dus is $p(x) = \lim
x = \mathrm{LIM}(x)$ daar; in het bijzonder $\mathrm{LIM} \leq p$
op $c$. Hahn--Banach (\cref{thm:b3:banach:hahnbanach}) zet
$\mathrm{LIM}$ voort tot $\ell^\infty_\R$ met overal
$\mathrm{LIM} \leq p$. Positiviteit: voor $x \geq 0$ is
$-\mathrm{LIM}(x) = \mathrm{LIM}(-x) \leq p(-x) =
\limsup\text{gem}(-x) \leq 0$. Verschuivingsinvariantie: de
middelen van $x - Sx$ telescoperen tot $\frac{x_1 - x_{n+1}}n \to
0$, dus $p(\pm(x - Sx)) = 0$ en $\mathrm{LIM}(x - Sx) = 0$.
Grenzen: $\mathrm{LIM}(x) \leq p(x) \leq \limsup x$ (de middelen
blijven achter bij de suprema), en dit op $-x$ toepassen geeft de
ondergrens.

\textbf{11.} $e_n \in c_0$, dus $\mathrm{LIM}(e_n) = 0$ voor alle
$n$. Was $\mathrm{LIM} = \Lambda_y$ met $y \in \ell^1$, dan was
$y_n = \Lambda_y(e_n) = 0$ voor alle $n$, dus $\Lambda_y = 0$;
maar $\mathrm{LIM}(\mathbf 1) = 1$. Dus $\mathrm{LIM} \in
(\ell^\infty)' \setminus \{\Lambda_y : y \in \ell^1\}$: opnieuw
$(\ell^\infty)' \neq \ell^1$.

\textbf{12.} Voor $x = (0,1,0,1,\dots)$ is $x + Sx = \mathbf 1$,
dus $2\,\mathrm{LIM}(x) = \mathrm{LIM}(x) + \mathrm{LIM}(Sx) = 1$
en $\mathrm{LIM}(x) = \frac12$. Was $\varphi$ een
verschuivingsinvariante multiplicatieve voortzetting van de
limiet, dan geeft $x\cdot Sx = 0$ dat $\varphi(x)\varphi(Sx) =
\varphi(x)^2 = 0$, dus $\varphi(x) = 0$; maar $\varphi(x) +
\varphi(Sx) = \varphi(\mathbf 1) = 1$ geeft $2\varphi(x) = 1$:
tegenspraak. Middelen en vermenigvuldigen gaan niet samen.

\textbf{13.} Zij $\norm{x^{(k)}}_p \leq M$. De coördinaten zijn
door $M$ begrensd: een diagonaalextractie geeft een deelrij (nog
steeds $x^{(k)}$ genoemd) met $x^{(k)}_n \to x_n$ voor elke $n$.
Dan is $x \in \ell^p$: $\sum_{n\leq N}\abs{x_n}^p =
\lim_k\sum_{n\leq N}\abs{x^{(k)}_n}^p \leq M^p$ voor alle $N$.
Zwakke convergentie: voor $y \in \ell^q$ en willekeurige $N$ is
\[
\abs{\Lambda_y(x^{(k)} - x)} \leq
\Bigl|\sum_{n \leq N}(x^{(k)}_n - x_n)y_n\Bigr|
+ 2M\Bigl(\sum_{n>N}\abs{y_n}^q\Bigr)^{1/q},
\]
waarbij de eerste term naar $0$ gaat als $k \to \infty$ (eindig
veel coördinaten) en de tweede klein is voor grote $N$ (Hölder op
de staart): dus $\Lambda_y(x^{(k)}) \to \Lambda_y(x)$ voor elke
$y$ --- zwakke convergentie, want elke functionaal is een
$\Lambda_y$ (vraag 6). In $\ell^1$ faalt dit: beschouw de
begrensde rij $(e_n)$. Elke deelrij $(e_{n_k})$ convergeert
coördinaatsgewijs naar $0$, dus is $0$ haar enige
kandidaat-limiet; maar toetsen tegen $y \in \ell^\infty$ met
$y_{n_k} = (-1)^k$ (en $0$ elders) geeft $\Lambda_y(e_{n_k}) =
(-1)^k$, wat divergeert. Er is dus geen zwak convergente deelrij:
zwakke \omterm{thm:b3:topology:metriccompact}{rijcompactheid} van de ballen kenmerkt de reflexieve
wereld.

\textbf{14.} Voor $y \in \ell^1$ is $\Lambda_y(x) = \sum x_ny_n$
gedefinieerd op $c_0$ met $\abs{\Lambda_y(x)} \leq \norm
x_\infty\norm y_1$, en toetsen op $x^{(N)} =
(\operatorname{sign}y_1, \dots, \operatorname{sign}y_N, 0, \dots)
\in c_0$ geeft $\Lambda_y(x^{(N)}) = \sum_{n\leq N}\abs{y_n} \to
\norm y_1$: dus $\norm{\Lambda_y} = \norm y_1$. Stel omgekeerd
$y_n = \Lambda(e_n)$ voor $\Lambda \in (c_0)'$; dezelfde toetsen
geven $\sum_{n \leq N}\abs{y_n} = \Lambda(x^{(N)}) \leq
\norm\Lambda$, dus $y \in \ell^1$, en $\Lambda = \Lambda_y$ op de
\omterm{def:b3:topology:interior}{dichte} eindige rijen en dus overal. De keten van \omterm{def:b3:banach:operator}{duale ruimten}:
$(c_0)' = \ell^1$, $(\ell^1)' = \ell^\infty$
(\cref{exo:b3:banach:10}), $(\ell^\infty)' \supsetneq \ell^1$
(vragen 9--11): de reflexiviteit faalt al bij de allereerste stap
--- $c_0'' = \ell^\infty \neq c_0$ --- en herstelt nooit meer.

\textbf{15.} $J x^{(k)} \in E'' = (E')'$ is een familie begrensde
functionalen op de \emph{banachruimte} $E'$ (\omterm{def:b3:banach:operator}{duale ruimten} zijn
\omterm{def:b3:complete:complete}{volledig}, \cref{prop:b3:banach:llcomplete}); en voor elke
$\Lambda \in E'$ convergeert de rij $Jx^{(k)}(\Lambda) =
\Lambda(x^{(k)})$ en is ze dus begrensd. Banach--Steinhaus op $E'$
geeft $\sup_k\norm{Jx^{(k)}} < \infty$, en $J$ is isometrisch
(\cref{rem:b3:banach:bidual}): dus $\sup_k\norm{x^{(k)}} <
\infty$.

\textbf{16.} ($\Rightarrow$) De begrensdheid is vraag 15; en de
coördinaten zijn de functionalen $\Lambda_{e_n}$.
($\Leftarrow$) Zij $M = \sup_k\norm{x^{(k)}}_p$, $y \in \ell^q$ en
$\varepsilon > 0$; kies $N$ met
$\bigl(\sum_{n>N}\abs{y_n}^q\bigr)^{1/q} < \varepsilon$. Dan is
\[
\abs{\Lambda_y(x^{(k)} - x)} \leq
\sum_{n\leq N}\abs{x^{(k)}_n - x_n}\,\abs{y_n} +
\norm{x^{(k)} - x}_p\,\varepsilon
\leq \sum_{n\leq N}\abs{x^{(k)}_n - x_n}\abs{y_n} +
(M + \norm x_p)\,\varepsilon,
\]
en de eindige som gaat naar $0$: dus $\limsup \leq (M + \norm
x_p)\varepsilon$ voor elke $\varepsilon$. (Dat $x \in \ell^p$ met
$\norm x_p \leq M$, volgt uit grenzen op eindige stukken in de
trant van Fatou: $\sum_{n \leq N}\abs{x_n}^p =
\lim_k\sum_{n\leq N}\abs{x^{(k)}_n}^p \leq M^p$.) Voor $e_k$ in
$\ell^2$: begrensd, coördinaatsgewijs $\to 0$, dus $e_k
\rightharpoonup 0$, en toch $\norm{e_k} = 1$: de eenheid massa
ontsnapt naar oneindige index, onzichtbaar voor elke vaste
functionaal.

\textbf{17.} Neem $\Lambda$ met $\norm\Lambda = 1$ en $\Lambda(x)
= \norm x$ (\cref{cor:b3:banach:hbnormed}). Dan is $\norm x =
\lim\Lambda(x^{(k)}) \leq \liminf\norm\Lambda\,\norm{x^{(k)}} =
\liminf\norm{x^{(k)}}$. (Met $e_k \rightharpoonup 0$: $0 \leq
\liminf 1$, en de ongelijkheid kan strikt zijn.)

\textbf{18.} In $\ell^2$ is $\norm{x^{(k)} - x}_2^2 =
\norm{x^{(k)}}_2^2 - 2\operatorname{Re}\langle x^{(k)}, x\rangle
+ \norm x_2^2$ (in het reële geval $-2\langle x^{(k)},
x\rangle$). Zwakke convergentie toegepast op de functionaal
$\Lambda_x$ geeft $\langle x^{(k)}, x\rangle \to \norm x_2^2$, en
de normen convergeren per hypothese: het rechterlid gaat dus naar
$\norm x^2 - 2\norm x^2 + \norm x^2 = 0$. Tegenvoorbeeld zonder
normconvergentie: $e_k \rightharpoonup 0$ met $\norm{e_k - 0} = 1
\not\to 0$.

\textbf{19.} Coördinaatsgewijze convergentie: pas de functionalen
$\Lambda_{e_n} \in (\ell^1)' = \ell^\infty$ toe (met $y = e_n$).
Constructie: zijn $k_{j-1}$ en $N_{j-1}$ gekozen, neem dan $k_j >
k_{j-1}$ zo groot dat $\sum_{n \leq N_{j-1}}\abs{x^{(k_j)}_n} <
\frac\delta{20}$ (eindig veel coördinaten, elk $\to 0$), en
daarna $N_j > N_{j-1}$ zo groot dat de staart voldoet aan
$\sum_{n > N_j}\abs{x^{(k_j)}_n} < \frac\delta{20}$ (convergentie
van de reeks die $\norm{x^{(k_j)}}_1$ definieert). Het blok $B_j =
\intoc{N_{j-1}}{N_j}$ draagt dan op $\frac\delta{10}$ na de hele
massa van $x^{(k_j)}$.

\textbf{20.} Met $y$ als gedefinieerd (overal $\abs{y_n} \leq 1$):
\[
\langle x^{(k_j)}, y\rangle
= \sum_{n \in B_j}\abs{x^{(k_j)}_n}
+ \sum_{n \notin B_j}x^{(k_j)}_ny_n
\geq \Bigl(\norm{x^{(k_j)}}_1 - \frac\delta{10}\Bigr) -
\frac\delta{10}
\geq \delta - \frac{2\delta}{10} = \frac{4\delta}5 > 0,
\]
waarbij de middelste ongelijkheid geldt omdat de massa buiten het
blok hoogstens $\frac\delta{10}$ is (vraag 19). Maar $y \in
\ell^\infty = (\ell^1)'$ en $x^{(k)} \rightharpoonup 0$ dwingen
$\langle x^{(k_j)}, y\rangle \to 0$ af: tegenspraak. Zwak naar nul
convergerende rijen van $\ell^1$ convergeren dus in norm naar nul,
en door translatie convergeren zwak convergente rijen in norm: de
stelling van Schur.

\textbf{21.} Een zwak convergente deelrij van $(e_k)$ zou limiet
$0$ hebben (coördinaten), dus volgens Schur $\norm{e_{k_j}}_1 \to
0$ --- maar de normen zijn $1$. Er is dus geen zwak convergente
deelrij, zoals in vraag 13 met de hand gevonden. Geen paradox:
Schur zegt dat rijen in $\ell^1$ de zwakke en de normtopologie
niet kunnen onderscheiden (de topologieën zelf verschillen wel
degelijk --- zwakke \omterm{def:b3:topology:topology}{omgevingen} zijn nooit normbegrensd), en de
zwakke \omterm{thm:b3:topology:metriccompact}{rijcompactheid} van de eenheidsbal is een andere, sterkere
eigenschap, gelijkwaardig met reflexiviteit
(Eberlein--Šmulian, buiten ons gereedschap; het falen hebben we
althans bewezen).

\textbf{22.} \emph{De inventaris.}
\begin{center}
\begin{tabular}{l|l|c|c|c}
$E$ & $E'$ & sep. & refl. & zwak rijcpt.\ ballen\\
\hline
$c_0$ & $\ell^1$ & ja & nee & nee ($e_1{+}\dots{+}e_k$)\\
$\ell^1$ & $\ell^\infty$ & ja & nee & nee ($e_k$, vr.~21)\\
$\ell^p$ & $\ell^q$ & ja & ja & ja (vr.~13)\\
$\ell^\infty$ & $\supsetneq\ell^1$ & nee & nee & nee\\
\end{tabular}
\end{center}
Kenmerken: $c_0$ --- haar biduale ruimte is $\ell^\infty$: de
eerste niet-reflexieve stap (vraag 14); $\ell^1$ --- de
eigenschap van Schur (vraag 20); $\ell^p$ --- reflexiviteit en
zwakke \omterm{def:b3:topology:compact}{compactheid} (vragen 7 en 13); $\ell^\infty$ ---
niet-separabel, met banachlimieten: functionalen die geen enkele
rij kan voorstellen (vragen 8 en 10--11). Eén familie ruimten,
vier verschillende werelden.

\textbf{23.} Zij $(f_k) \subseteq F$ met $\norm{x - f_k} \to d$.
Dan is $\norm{f_k} \leq \norm x + \sup_k\norm{x - f_k}$:
begrensd, dus levert vraag 13 een deelrij $f_{k_j}
\rightharpoonup f$. Was $f \notin F$, dan is $\delta =
\operatorname{dist}(f, F) > 0$ ($F$ is gesloten); op $F \oplus \R
f$ voldoet de lineaire vorm $\lambda(g + tf) = t$ aan
$\abs{\lambda(u)} \leq \norm u/\delta$ (want $\norm{g + tf} \geq
\abs t\,\delta$), en Hahn--Banach zet die voort tot $\Lambda \in
(\ell^p)'$ met $\Lambda\restriction_F = 0$ en $\Lambda(f) = 1$;
maar dan is $0 = \Lambda(f_{k_j}) \to \Lambda(f) = 1$:
tegenspraak. Dus $f \in F$, en $x - f_{k_j} \rightharpoonup x - f$
geeft met vraag 17
\[
d \leq \norm{x - f} \leq \liminf_j\,\norm{x - f_{k_j}} = d :
\]
de afstand wordt in $f$ bereikt. In $c_0$: $\abs{\Lambda(x)} \leq
\sum_n2^{-n}\abs{x_n} < \norm x_\infty$ voor elke $x \neq 0$ (een
nulrij $\neq 0$ kan niet voor alle $n$ aan $\abs{x_n} = \norm
x_\infty$ voldoen), terwijl de afgeknotte rijen $(1, \dots, 1, 0,
\dots)$ geven $\Lambda = 1 - 2^{-N} \to 1$: dus $\norm\Lambda =
1$, nooit bereikt. Afstandsformule: voor $f \in \ker\Lambda$ is
$\abs{\Lambda(x)} = \abs{\Lambda(x - f)} \leq \norm{x - f}$, dus
$\operatorname{dist}(x, \ker\Lambda) \geq \abs{\Lambda(x)}$; en
omgekeerd ligt voor $u$ in de eenheidsbal met $\Lambda(u) \geq 1 -
\varepsilon$ de vector $f = x - \frac{\Lambda(x)}{\Lambda(u)}u$ in
$\ker\Lambda$ met $\norm{x - f} \leq \frac{\abs{\Lambda(x)}}{1 -
\varepsilon}$: gelijkheid. Werd ze in een $f \in \ker\Lambda$
bereikt, dan voldeed $z = x - f$ aan $\abs{\Lambda(z)} =
\abs{\Lambda(x)} = \norm z \neq 0$, zodat $\Lambda$ haar norm in
$z/\norm z$ zou bereiken: onmogelijk. Een gesloten hypervlak van
$c_0$ zonder ergens een dichtstbijzijnd punt --- de reflexiviteit
was geen versiering.

\textbf{24.} Stel $y_1 = x^{(1)}$. Gegeven $y_1, \dots, y_j$ gaat
elke afbeelding $k \mapsto \langle y_i, x^{(k)}\rangle$ naar $0$
($y_i \in \ell^2 = (\ell^2)'$), dus is er voorbij de vorige index
een $k_{j+1}$ met $\abs{\langle y_i, x^{(k_{j+1})}\rangle} \leq
\frac1{j+1}$ voor $i = 1, \dots, j$; noem die keuze $y_{j+1}$. Dan
is
\[
\Bigl\lVert\sum_{j=1}^my_j\Bigr\rVert_2^2
= \sum_{j=1}^m\norm{y_j}_2^2
+ 2\sum_{j=2}^m\sum_{i<j}\langle y_i, y_j\rangle
\leq mC^2 + 2\sum_{j=2}^m\frac{j-1}j
\leq mC^2 + 2m,
\]
en delen door $m^2$ geeft $\norm{\frac1m\sum_jy_j}_2^2 \leq
\frac{C^2 + 2}m \to 0$. (Voor een zwakke limiet $x \neq 0$: pas
dit toe op $x^{(k)} - x$.) Op de orthonormale $(e_k)$ is zelfs
geen extractie nodig: $\norm{\frac1m(e_1 + \dots + e_m)}_2 =
\frac{\sqrt m}m = \frac1{\sqrt m}$. Middelen zet zwakke
convergentie dus om in normconvergentie: de eigenschap van
Banach--Saks van $\ell^2$.

\textbf{25.} $A_mx = \frac1m\sum_{i=0}^{m-1}S^ix$, dus geven de
lineariteit en de verschuivingsinvariantie $L(A_mx) = L(x)$. Kies
voor een begrensde $u$ en $\varepsilon > 0$ een $N$ met $u_n \leq
\limsup u + \varepsilon$ voor $n \geq N$; de positiviteit
toegepast op $(\limsup u + \varepsilon)\mathbf 1 - S^Nu \geq 0$
samen met $L(S^Nu) = L(u)$ geeft $L(u) \leq \limsup u +
\varepsilon$, en symmetrisch $L(u) \geq \liminf u - \varepsilon$:
dus $\liminf A_mx \leq L(x) \leq \limsup A_mx$ voor elke $m$. Is
$x$ $T$-periodiek, dan is $A_Tx$ de constante rij gelijk aan het
periodegemiddelde $\mu$: dus $L(x) = \mu$ voor \emph{elke}
banachlimiet --- $\frac13$ op $(1,0,0,\dots)$ en $\frac12$ op
$(0,1,0,1,\dots)$ als in vraag 12. Voor de blokrij: bij $N = 3^j$
met $j$ even bestaat het laatste blok geheel uit enen, dus is het
cesàromiddel $\geq \frac{3^j - 3^{j-1}}{3^j} = \frac23$; en bij $N
= 3^j$ met $j$ oneven liggen alle enen in $\intoc0{3^{j-1}}$, dus
is het middel $\leq \frac13$. Bijgevolg is $p(x) \geq \frac23$ en
$-p(-x) = \liminf_n\frac{x_1 + \dots + x_n}n \leq \frac13$.
Definieer op $M = c \oplus \R x$: $\Lambda_+(y + tx) = \lim y +
t\,p(x)$. Dominatie door $p$: voor $t > 0$ geeft de sublineariteit
$p(tx) \leq p(y + tx) + p(-y)$, oftewel $p(y + tx) \geq t\,p(x) +
\lim y$ (merk op dat $p(\pm y) = \pm\lim y$ voor $y \in c$: de
cesàromiddelen van een convergente rij convergeren naar haar
limiet); voor $t = -s < 0$ geeft $p(y) \leq p(y - sx) + p(sx)$ dat
$p(y - sx) \geq \lim y - s\,p(x)$; en voor $t = 0$ is er
gelijkheid. Dus is $\Lambda_+ \leq p$ op $M$, en Hahn--Banach zet
haar voort tot $L_+ \leq p$ op $\ell^\infty_\R$, wat precies als
in vraag 10 een banachlimiet is (uit de dominatie door $p$ volgen
de positiviteit, de verschuivingsinvariantie en de waarde $\lim$
op $c$), met $L_+(x) = p(x) \geq \frac23$. Dezelfde berekening met
$\Lambda_-(y + tx) = \lim y - t\,p(-x)$ (met $p(-sx) \leq p(y -
sx) + p(-y)$ voor het geval $t = -s < 0$) levert een banachlimiet
$L_-$ met $L_-(x) = -p(-x) \leq \frac13$. Twee banachlimieten, één
rij, twee waarden: buiten de periodieke (en algemener de
\emph{bijna convergente}) wereld is een banachlimiet een echte
keuze.
\end{solution}
